Un reactor idealizado contiene \(1\,\mathrm L\) de mezcla a temperatura constante. Antes de dosificar hay \(3\,\mathrm{mol}\) de A y no hay B ni P. La única reacción es \(\mathrm{A+B\longrightarrow P}\), con estequiometría uno a uno. La velocidad volumétrica de formación de P es \(k c_Ac_B\), donde \(c_A,c_B\) son las concentraciones y \(k=1\,\mathrm{L/(mol\,min)}\). Los datos son sintéticos.

Se comparan dos programas que añaden en total \(2\,\mathrm{mol}\) de B:
\[
\begin{array}{c|cc}
\text{programa}&t=0&t=0.5\,\mathrm{min}\\\hline
\mathrm I&2\,\mathrm{mol}&0\\
\mathrm{II}&1\,\mathrm{mol}&1\,\mathrm{mol}
\end{array}
\]
Las adiciones se idealizan como instantáneas, de volumen despreciable y mezcla inmediata. Durante el instante de adición sólo cambia la cantidad de B por la dosis indicada; la reacción avanza entre adiciones. No hay salida ni otras reacciones.

Se requiere obtener al menos \(1.40\,\mathrm{mol}\) de P al cabo de un minuto y mantener la cantidad de B libre por debajo o igual a \(1.50\,\mathrm{mol}\) en toda la operación, también inmediatamente después de cada carga. Formule los balances y las condiciones de unión, compare el producto obtenido y el máximo de B para cada programa, y decida cuáles cumplen los criterios. Explique por qué la fecha de evaluación importa al comparar dosis totales iguales. Puede usar calculadora.
